% ==============================
% Chapter 2 – Problem Definition
% ==============================

\chapter{Problem Definition}

% --------------------------------
% 2.1 Problem Statement
% --------------------------------

\section{Problem Statement}

The central barrier holding back quantum software development today is the gap between abstract quantum algorithms and physical hardware constraints. When researchers publish an algorithm like Grover's search or Quantum Phase Estimation, they describe circuits using theoretical gates acting on arbitrary pairs of qubits. In reality, physical quantum processors cannot execute arbitrary multi-qubit interactions.

On superconducting chips, such as those fabricated by IBM and Google, qubits exist as physical lithographic circuits on a wafer. Qubit interactions rely on fixed or tunable microwave resonators that couple only physically adjacent qubits. If an algorithm requires a two-qubit gate (such as CNOT or CZ) between qubit $q_0$ and qubit $q_3$ on a linear chip layout ($q_0 - q_1 - q_2 - q_3$), the hardware cannot physically execute the gate. The physical transmission lines simply do not connect them.

To run this circuit, someone must route the quantum state of $q_0$ across the chip. This requires inserting a cascade of SWAP gates: swapping $q_0$ with $q_1$, then $q_1$ with $q_2$, executing the interaction, and swapping back. Each SWAP operation is costly; on native hardware, a single SWAP decomposes into three consecutive CNOT gates ($CX_{01} CX_{10} CX_{01}$). On NISQ-era processors with two-qubit gate error rates ranging between 0.5\% and 3.0\%, adding just two or three SWAP gates quickly doubles or triples the circuit depth and swamps the true quantum signal with physical decoherence noise.

Currently, developers face two frustrating choices:
\begin{enumerate}[leftmargin=*]
    \item Manually re-route circuits by hand for every target chip architecture, which is tedious, prone to indexing mistakes, and fragile across chip revisions.
    \item Rely on complex, heavyweight commercial toolkits that obscure compiler decisions, require cloud credentials, or tie the developer directly into proprietary hardware ecosystems.
\end{enumerate}

There is an acute lack of transparent, lightweight, and hardware-agnostic execution runtimes that allow students, researchers, and software engineers to compile and inspect quantum circuits locally without enterprise overhead or cloud lock-in.

% --------------------------------
% 2.2 Proposed Solution Overview
% --------------------------------

\section{Proposed Solution Overview}

To solve this problem, we designed and implemented the \textbf{Quantum Virtual Machine (QVM)}. QVM acts as an intelligent middleware layer positioned between high-level quantum circuit descriptions and arbitrary hardware execution targets.

Instead of writing hardware-specific scripts, developers define their quantum logic using standard formats: OpenQASM 3.0 \cite{cross2022qasm3}, OpenQASM 2.0 \cite{cross2017qasm}, or JSON. QVM processes the circuit through a structured three-stage pipeline:

\begin{enumerate}[leftmargin=*]
    \item \textbf{Ingestion and Normalization:} The parser ingests the input file, validates its grammar, and builds an in-memory Abstract Syntax Tree (AST). It converts declarations and operations into a standardized \texttt{QuantumCircuit} Intermediate Representation (IR). Non-primitive operations (such as Toffoli or multi-controlled gates) are decomposed into target-compatible basis gates.
    \item \textbf{Topology-Aware Transpilation:} Given a target hardware coupling graph (such as a 5-qubit linear chain, a 2D grid, or custom coupling rules), the transpiler identifies two-qubit operations that violate adjacency constraints. It computes shortest routing paths, evaluates lookahead heuristics (Greedy BFS or SABRE \cite{li2019sabre}), inserts the minimal necessary SWAP gates, and tracks the evolving mapping between logical and physical qubits.
    \item \textbf{Verification and Execution:} Once compiled, the physical circuit can either be exported back to OpenQASM for external deployment or executed locally on QVM's dual simulation engines: an exact statevector simulator for up to 12 qubits, or an MPS tensor network simulator for larger systems with low entanglement.
\end{enumerate}

This pipeline decouples the algorithmic logic from the execution backend, realizing true ``Write Once, Run Anywhere'' portability.

% --------------------------------
% 2.3 Objectives of the Proposed System
% --------------------------------

\section{Objectives of the Proposed System}

The design and development of QVM was guided by five core engineering objectives:

\begin{itemize}[leftmargin=*]
    \item \textbf{BO-1: Hardware-Agnostic Compilation:} Enable developers to write quantum programs once in standard OpenQASM or JSON and automatically compile them to distinct hardware topologies (Linear, 2D Grid, Ring) without modifying the original circuit logic.
    \item \textbf{BO-2: Compiler Transparency and Pedagogical Clarity:} Expose clear diagnostics and step-by-step intermediate representations showing exactly how logical qubits map to physical qubits, which decomposition rules fired, and where SWAP gates were inserted.
    \item \textbf{BO-3: Lightweight, Standalone Local Execution:} Provide a fully self-contained runtime under 50~MB that runs on commodity personal computers with zero mandatory cloud accounts, API tokens, or external hardware subscriptions.
    \item \textbf{BO-4: Standard Interoperability:} Adhere strictly to OpenQASM 2.0 and 3.0 specifications, including classical bit registers, mid-circuit measurement collapse, and conditional branching, ensuring compiled output can be exported directly to other quantum tools.
    \item \textbf{BO-5: Automated Gate Decomposition and Optimization:} Automate the reduction of high-level unitary operations into elementary one- and two-qubit basis gates, incorporating lookahead heuristics to minimize total circuit depth and SWAP overhead.
\end{itemize}

% --------------------------------
% 2.4 Scope of the System
% --------------------------------

\section{Scope of the System}

Setting explicit project boundaries kept our development focused on numerical accuracy and compilation correctness within the nine-month Final Year Project timeline.

\textbf{What is in Scope:}
\begin{itemize}[leftmargin=*]
    \item \textbf{Input Ingestion:} Full parsing of OpenQASM 2.0, OpenQASM 3.0 (including classical variables, \texttt{if} conditionals, and loops), and JSON gate lists.
    \item \textbf{Circuit Manipulation:} Gate decomposition (Toffoli, controlled-$U$, multi-controlled gates) into elementary gates ($\{H, X, Y, Z, Rx, Ry, Rz, CX, CZ, SWAP\}$).
    \item \textbf{Transpilation:} Routing across user-defined or pre-set hardware coupling graphs using Greedy BFS and SABRE heuristic lookahead.
    \item \textbf{Dual Simulation Backends:} Exact statevector simulation for circuits up to 12 qubits (with classical memory registers and mid-circuit collapse) and Matrix Product State (MPS) tensor network simulation for circuits scaling to 20+ qubits when entanglement is bounded.
    \item \textbf{Noise Simulation:} Kraus operator channels for depolarizing noise, $T_1$ relaxation, $T_2$ dephasing, and asymmetric readout errors.
    \item \textbf{Interfaces:} A command-line tool, an asynchronous FastAPI REST service, and an interactive Next.js web dashboard.
\end{itemize}

\textbf{What is out of Scope:}
\begin{itemize}[leftmargin=*]
    \item \textbf{Pulse-Level Hardware Control:} QVM operates at the logical and gate-level circuit abstraction. It does not synthesize low-level microwave analog pulse envelopes or frequency-shift calibrations.
    \item \textbf{Active Fault-Tolerant Quantum Error Correction (QEC):} The system does not implement live real-time surface-code syndrome decoding loops.
    \item \textbf{Direct Proprietary Hardware Submissions:} While QVM exports valid OpenQASM files ready for submission to IBM Quantum or AWS Braket, it intentionally avoids managing remote cloud job queues or enterprise billing credentials.
\end{itemize}

% --------------------------------
% 2.5 System Architecture and Modules
% --------------------------------

\section{System Architecture and Modules}

The Quantum Virtual Machine is structured as a pipeline of independent, loosely coupled modules. Each module performs a single well-defined transformation on the quantum program:

\begin{enumerate}[leftmargin=*]
    \item \textbf{Input & Interface Layer:} Ingests program files via CLI, REST API requests, or web editor interactions.
    \item \textbf{Parser Module:} Converts raw textual descriptions into a validated Abstract Syntax Tree using Lark formal grammars.
    \item \textbf{Intermediate Representation (IR):} Normalizes the AST into a \texttt{QuantumCircuit} instance managing gate streams, qubit declarations, and classical memory.
    \item \textbf{Decomposer Module:} Synthesizes multi-qubit and high-level gates into native elementary primitives.
    \item \textbf{Transpiler Module:} Resolves topological coupling constraints by mapping logical qubits and inserting SWAP gates via graph algorithms.
    \item \textbf{Simulation Engines:} Executes physical circuits using either exact statevector array mathematics or tensor-network SVD contractions.
    \item \textbf{Output & Visualization Layer:} Computes measurement statistics, renders circuit schematics, and formats JSON/ASCII responses.
\end{enumerate}

% --------------------------------
% 2.5.1 Module Descriptions
% --------------------------------

\subsection{Module 1: Parser Module}

The Parser Module handles all input ingestion, serving as the gateway to the virtual machine. It accepts circuits written in OpenQASM 2.0 \cite{cross2017qasm}, OpenQASM 3.0 \cite{cross2022qasm3}, or structured JSON.

For OpenQASM 3.0, the module relies on the Lark parsing toolkit \cite{lark2023} configured with an EBNF grammar file (\texttt{qasm3.lark}). The parser executes in two passes: the first pass discovers register declarations (qubits and classical bit arrays), while the second pass walks the AST to parse gates, measurement assignments, conditional branches, and loops. If a user provides malformed syntax or references undeclared qubits, the parser catches the error immediately and outputs line numbers with contextual diagnostic suggestions.

\textbf{Functional Requirements:}
\begin{itemize}[leftmargin=*]
    \item \textbf{FR-1.1:} The system shall parse OpenQASM 2.0 circuits containing standard single- and two-qubit gates, barriers, and measurement directives.
    \item \textbf{FR-1.2:} The system shall parse OpenQASM 3.0 programs containing classical registers, \texttt{if} conditionals, and loop constructs.
    \item \textbf{FR-1.3:} The system shall validate circuit syntax and provide descriptive error messages indicating exact line and column numbers upon failure.
    \item \textbf{FR-1.4:} The system shall construct a standardized \texttt{QuantumCircuit} Intermediate Representation from valid parsed inputs.
    \item \textbf{FR-1.5:} The system shall accept JSON gate list representations for testing and automated tool integration.
\end{itemize}

\subsection{Module 2: Transpiler Module}

The Transpiler Module adapts hardware-agnostic circuits to the physical connectivity limitations of target quantum processors \cite{cowtan2019}. It takes a logical circuit and a hardware coupling graph as inputs.

If an operation requires a two-qubit gate between qubits that share no direct physical edge, the transpiler intervenes. It maintains two dynamic mapping tables: forward (logical $\to$ physical) and inverse (physical $\to$ logical). Users can pick between two routing algorithms:
\begin{enumerate}
    \item \textbf{Greedy BFS:} Computes the shortest path across the hardware graph using Breadth-First Search and inserts SWAP gates along the path. It runs in milliseconds and is ideal for fast local debugging.
    \item \textbf{SABRE Routing:} Evaluates future gates in a lookahead queue using a decay-weighted cost heuristic \cite{li2019sabre}. By anticipating upcoming multi-qubit interactions, SABRE reduces unnecessary SWAPs by 30\% to 40\% compared to the greedy baseline.
\end{enumerate}

\textbf{Functional Requirements:}
\begin{itemize}[leftmargin=*]
    \item \textbf{FR-2.1:} The system shall accept hardware coupling graph specifications defining valid adjacent qubit pairs.
    \item \textbf{FR-2.2:} The system shall implement Greedy BFS routing to route non-adjacent gates via shortest paths.
    \item \textbf{FR-2.3:} The system shall implement the SABRE lookahead heuristic to minimize total SWAP gate count and circuit depth.
    \item \textbf{FR-2.4:} The system shall allow users to select their preferred routing strategy and configure lookahead parameters.
    \item \textbf{FR-2.5:} The system shall provide an optional mapping restoration pass that returns physical qubits to their original logical indices at circuit completion.
\end{itemize}

\subsection{Module 3: Simulator Module}

The Simulator Module provides execution backends to verify and study quantum circuit behavior. Rather than relying on a single engine, QVM incorporates dual simulation approaches:

\begin{itemize}[leftmargin=*]
    \item \textbf{Statevector Simulator:} Represents the complete quantum state as a complex NumPy array of size $2^n$ \cite{numpy2020}. It optimizes gate execution through multidimensional tensor reshaping and axis permutations, avoiding expensive $2^n \times 2^n$ matrix allocations. It supports projective mid-circuit measurements according to the Born rule, collapsing the statevector and updating classical registers in real time.
    \item \textbf{Matrix Product State (MPS) Simulator:} Models the quantum state as a 1D tensor train \cite{vidal2003mps}. When a two-qubit gate fires, the engine contracts adjacent rank-3 tensors and performs Singular Value Decomposition (SVD), discarding singular values below a threshold to cap bond dimension $\chi$. This allows simulating 20+ qubits on desktop machines as long as entanglement remains localized.
    \item \textbf{Kraus Noise Engine:} Applies open-system noise channels (depolarizing, amplitude damping, phase damping, and readout classification confusion) to evaluate how physical imperfections degrade circuit outputs.
\end{itemize}

\textbf{Functional Requirements:}
\begin{itemize}[leftmargin=*]
    \item \textbf{FR-3.1:} The system shall perform exact statevector simulation for circuits up to 12 qubits on standard desktop hardware.
    \item \textbf{FR-3.2:} The system shall perform Matrix Product State (MPS) simulation with SVD truncation for low-entanglement circuits up to 20+ qubits.
    \item \textbf{FR-3.3:} The system shall execute mid-circuit measurements with statevector collapse and write outcomes to classical memory registers.
    \item \textbf{FR-3.4:} The system shall evaluate classical conditional statements during simulation to support hybrid feedback algorithms.
    \item \textbf{FR-3.5:} The system shall compute analytic state probabilities and simulate stochastic shot-based sampling with configurable seeds.
\end{itemize}

\subsection{Module 4: Visualization Module}

The Visualization Module translates raw numerical statevectors and circuit data structures into intuitive graphical diagrams using Matplotlib \cite{matplotlib2007}.

The module provides two primary visual representations:
\begin{enumerate}
    \item \textbf{Circuit Diagram Renderer:} Plots horizontal qubit wires, labeled gate blocks, multi-qubit control points, target symbols, and classical measurement wires.
    \item \textbf{Measurement Histogram Plotter:} Displays probability distributions across basis states, filtering out near-zero amplitudes to keep charts readable.
\end{enumerate}

\textbf{Functional Requirements:}
\begin{itemize}[leftmargin=*]
    \item \textbf{FR-4.1:} The system shall render graphical circuit diagrams displaying qubit wires, gate symbols, and measurement operations.
    \item \textbf{FR-4.2:} The system shall generate probability histograms illustrating measurement outcome distributions.
    \item \textbf{FR-4.3:} The system shall support customizable plot dimensions, color schemes, and output formats (PNG, SVG, PDF).
    \item \textbf{FR-4.4:} The system shall export rendered visual artifacts to disk for inclusion in documentation and academic reports.
\end{itemize}

\subsection{Module 5: CLI and Web API Module}

This module delivers user interfaces suited for diverse operational needs: automated scripting, web integration, and interactive browser experimentation.

The CLI (\texttt{src/qvm/cli.py}) uses Python's \texttt{argparse} to expose compilation flags, routing choices, shot counts, and noise parameters directly from the terminal. The Web API (\texttt{api/app.py}), developed with FastAPI \cite{fastapi2023}, exposes asynchronous HTTP endpoints (\texttt{POST /run}, \texttt{GET /health}, \texttt{GET /circuits}) with automatic JSON schema validation via Pydantic. It connects with a Supabase cloud database to persist user-submitted circuits and query past execution traces.

\textbf{Functional Requirements:}
\begin{itemize}[leftmargin=*]
    \item \textbf{FR-5.1:} The system shall provide a command-line interface supporting file arguments, routing flags, and visualization options.
    \item \textbf{FR-5.2:} The system shall implement a RESTful web API accepting JSON circuit payloads and returning execution results.
    \item \textbf{FR-5.3:} The system shall generate interactive OpenAPI documentation for all exposed HTTP endpoints.
    \item \textbf{FR-5.4:} The system shall provide an interactive web dashboard for circuit editing, preset loading, and result visualization.
    \item \textbf{FR-5.5:} The system shall execute simulation requests asynchronously without blocking concurrent server operations.
\end{itemize}

% --------------------------------
% 2.7 Assumptions and Dependencies
% --------------------------------

\section{Assumptions and Dependencies}

To ensure realistic expectations during deployment and evaluation, we identified the following operating assumptions and technical dependencies:

\textbf{Assumptions:}
\begin{itemize}[leftmargin=*]
    \item Users have a foundational understanding of quantum concepts (superposition, measurement, basic gate operations) and basic familiarity with the command line or web browsers.
    \item Input circuits specify mathematically valid unitary gates and non-negative qubit indices.
    \item The host machine provides at least 4~GB of available RAM for simulating circuits up to 10 qubits, and 8~GB for 12-qubit statevectors.
\end{itemize}

\textbf{Dependencies:}
\begin{itemize}[leftmargin=*]
    \item \textbf{Runtime Environment:} Python 3.10 or higher.
    \item \textbf{Core Libraries:} NumPy $\ge$ 1.24 (linear algebra), Lark $\ge$ 1.1 (parsing), FastAPI $\ge$ 0.100 (API service), Matplotlib $\ge$ 3.7 (plotting), and pytest $\ge$ 7.4 (testing framework).
    \item \textbf{Standards Adherence:} Adherence to OpenQASM 2.0 and 3.0 specification documents for grammar compliance and circuit interchange.
\end{itemize}

% --------------------------------
% 2.8 Chapter Summary
% --------------------------------

\section{Chapter Summary}

This chapter laid out the practical bottlenecks we tackle in the Quantum Virtual Machine: bridging the divide between abstract gate sequences and physical coupling limits. We examined why planar coupling graphs force compilers to insert SWAP cascades, and how unoptimized routing compounds gate errors on NISQ devices. We then mapped out QVM's three-stage compilation pipeline, set five verifiable project objectives, scoped the system boundaries, and indexed twenty-five functional requirements across five core subsystems. With these operational constraints established, Chapter 3 formalizes the functional requirements, user classes, and use case model.
